\chapter{有限脉冲：从单色轨迹相位到延迟依赖边带}
\label{chap:pulses}

有限脉冲不是新的电子动力学。第~\ref{chap:pinem-core} 章已经给出任意时空外场的一阶特征线解；这里唯一的新输入是光场具有慢包络。为了避免重复推导一阶输运方程，先把随动坐标直接代入已经得到的轨迹积分，再逐项说明准单色近似、$A$--$E$ 转换、穿越时间近似和边带可分辨条件。

固定一条纵向特征线，定义
\begin{equation}
\zeta\equiv z-v_0t,\qquad z(t')=\zeta+v_0t'.
\label{eq:moving-coordinate}
\end{equation}
在时间规范 $\phi=0$ 下，式~\eqref{eq:eikonal-phase-general} 直接给出
\begin{align}
\frac{g(\zeta,t)}{g(\zeta,t_0)}
&=\exp\left[
\frac{\ii q_{\rm e}v_0}{\hbar}
\int_{t_0}^{t} A_z(\zeta+v_0t',t')\,\dd t'
\right].
\label{eq:g-integral}
\end{align}
这里没有重新建立运动方程：只是把第~\ref{chap:pinem-core} 章的通用特征线写成一维随动坐标。下面所有脉冲修正都从这个已经闭合的母式出发。

\section{把慢光学包络代入同一条轨迹积分}

对准单色脉冲，必须区分“频域 phasor”与“含慢包络的正频时域场”。把前文的空间 phasor 记作 $E_{\omega,z}(z)$，并定义正频部分
\begin{equation}
 E_z^{(+)}(z,t)
 \equiv E_{\omega,z}(z)\,
 a_{\rm L}(t-\tau)\,
 \ee^{-\ii\omega t},
 \qquad
 a_{\rm L}(t-\tau)
 =\exp\left[-\frac{(t-\tau)^2}{4\sigma_{\mathrm L}^2}\right].
 \label{eq:complex-field}
\end{equation}
真实电场仍是 $E_z=\operatorname{Re}E_z^{(+)}$。准单色条件不是一句“脉冲很长”，而是慢包络的相对变化率小于载频：
\begin{equation}
 \boxed{
 \epsilon_{\rm SVEA}(t)
 \equiv
 \frac{|\partial_t a_{\rm L}(t-\tau)|}
 {\omega |a_{\rm L}(t-\tau)|}
 =
 \frac{|t-\tau|}{2\omega\sigma_{\rm L}^2}
 \ll1}
 \label{eq:SVEA-field-condition}
\end{equation}
在脉冲主要支撑区 $|t-\tau|\lesssim\sigma_{\rm L}$ 内，一个简单充分条件就是
\begin{equation}
 \boxed{\omega\sigma_{\rm L}\gg1.}
 \label{eq:quasimonochromatic-condition}
\end{equation}

但式~\eqref{eq:complex-field} 还隐含了另一件与“载频远大于带宽”不同的事情：整个脉冲带宽内，纳米结构的空间场型也必须近似不变，否则不能把完整频谱写成一个固定的 $E_{\omega,z}(z)$ 乘同一个时间包络。记脉冲特征带宽为 $\Delta\omega_{\rm L}\sim\sigma_{\rm L}^{-1}$，并用 $\|\cdot\|_{\rm int}$ 表示在电子实际采样区域上的一致场范数，则一个局部充分条件是
\begin{equation}
\boxed{
\epsilon_{\rm mode}
\equiv
\Delta\omega_{\rm L}
\frac{\|\partial_\omega E_{\omega,z}\|_{\rm int}}
{\|E_{\omega,z}\|_{\rm int}}
\ll1.}
\label{eq:pulse-spatial-mode-freeze}
\end{equation}
式~\eqref{eq:quasimonochromatic-condition} 控制时间载波与慢包络的分离，式~\eqref{eq:pulse-spatial-mode-freeze} 控制材料色散和几何共振在脉冲带宽内是否可以冻结。靠近窄线宽 Mie 或材料共振时，前者可能成立而后者失效；此时必须保留 $\bm E_{\Omega}(\bm r)$ 的频率依赖并对光谱逐频积分，而不能只给单频 $\beta$ 乘一个时间 Gaussian。

为了把一阶输运式~\eqref{eq:g-integral} 写成电场形式，这里选择 $\phi=0$ 的时间规范；这是计算上的规范选择，不是新的物理近似。相应正频矢势写成
\begin{equation}
 A_z^{(+)}(z,t)
 =\mathcal A_z(z,t)\ee^{-\ii\omega t}.
\end{equation}
由 $E_z^{(+)}=-\partial_tA_z^{(+)}$，
\begin{align}
E_z^{(+)}
&=-\left[(\partial_t\mathcal A_z)-\ii\omega\mathcal A_z\right]
 \ee^{-\ii\omega t}\\
&=\left[\ii\omega\mathcal A_z-\partial_t\mathcal A_z\right]
 \ee^{-\ii\omega t}.
\end{align}
在式~\eqref{eq:SVEA-field-condition} 控制下，$\partial_t\mathcal A_z$ 相对 $\omega\mathcal A_z$ 是一阶小量。为了不把无量纲误差直接加到有量纲矢势上，显式定义余项 $\delta A_z^{(+)}$：
\begin{equation}
 A_z^{(+)}(z,t)
 =\frac{E_z^{(+)}(z,t)}{\ii\omega}
 +\delta A_z^{(+)}(z,t),
 \qquad
 \frac{|\delta A_z^{(+)}|}{|E_z^{(+)}|/\omega}
 =O(\epsilon_{\rm SVEA}).
 \label{eq:A-from-E-positive-frequency}
\end{equation}
取实部得到真正代入电子相位积分的实矢势
\begin{equation}
 A_z(z,t)
 =\frac{1}{\omega}\operatorname{Im}E_z^{(+)}(z,t)
 +\delta A_z(z,t),
 \qquad
 |\delta A_z|=O\!\left(\epsilon_{\rm SVEA}\frac{|E_z^{(+)}|}{\omega}\right).
 \label{eq:A-from-E}
\end{equation}
因此 SVEA 的误差不是一个孤立的无量纲加数，而是相对于主矢势 $|E^{(+)}|/\omega$ 的相对误差。

将式~\eqref{eq:complex-field} 与 \eqref{eq:A-from-E} 逐项代入式~\eqref{eq:g-integral}：
\begin{align}
 \frac{g(z',t)}{g(z',t_0)}
 =\exp\Bigg\{\ii\frac{q_{\rm e}v_0}{\hbar\omega}
 \int_{t_0}^{t}\dd t'\,
 \operatorname{Im}\Bigg[
 &E_{\omega,z}(z'+v_0t')\ee^{-\ii\omega t'}\notag\\
 &\times\exp\left(-\frac{(t'-\tau)^2}{4\sigma_{\mathrm L}^2}\right)
 \Bigg]\Bigg\}.
 \label{eq:time-integral-E}
\end{align}
令
\begin{equation}
 z''=z'+v_0t',
 \qquad
 t'=\frac{z''-z'}{v_0},
 \qquad
 \dd t'=\frac{\dd z''}{v_0}.
 \label{eq:zpp-change}
\end{equation}
则
\begin{align}
 \ee^{-\ii\omega t'}
 &=\exp\left[-\ii\frac{\omega}{v_0}(z''-z')\right]\\
 &=\ee^{+\ii(\omega/v_0)z'}
 \ee^{-\ii(\omega/v_0)z''},
 \label{eq:phase-change}
\end{align}
而脉冲包络变成
\begin{equation}
 \exp\left[-\frac{(z''-z'-v_0\tau)^2}{4v_0^2\sigma_{\mathrm L}^2}\right].
 \label{eq:pulse-in-zpp}
\end{equation}
到这一步还没有假定光脉冲在电子穿越近场期间保持不变。先把有限积分区间扩展到整条轨迹。若初、末时刻选在相互作用完全结束之后，使截去的尾部满足
\begin{equation}
 \boxed{
 \int_{-\infty}^{t_0}\!|A_z(z'+v_0t',t')|\,\dd t'
 +\int_t^{+\infty}\!|A_z(z'+v_0t',t')|\,\dd t'
 \ll
 \int_{t_0}^{t}\!|A_z(z'+v_0t',t')|\,\dd t',}
 \label{eq:infinite-limit-condition}
\end{equation}
就可以令 $t_0\to-\infty$、$t\to+\infty$。把式~\eqref{eq:phase-change} 和脉冲包络~\eqref{eq:pulse-in-zpp} 代回式~\eqref{eq:time-integral-E}。为了同时记录 $A$--$E$ 转换带来的误差，定义相互作用区上的最大慢变参数与主相位尺度
\begin{align}
\epsilon_{\rm SVEA}^{\max}
&\equiv \sup_{t'\in I_{\rm int}}\epsilon_{\rm SVEA}(t'),\\
\Phi_*
&\equiv
\frac{|q_{\rm e}|v_0}{\hbar\omega}
\int_{I_{\rm int}}\!\dd t'\,
|E_z^{(+)}(z'+v_0t',t')|,
\label{eq:pulse-phase-error-scale}
\end{align}
其中 $I_{\rm int}$ 是场与电子显著重叠的时间区间。由式~\eqref{eq:A-from-E}，遗漏矢势在指数中造成的无量纲相位误差至多为 $O(\epsilon_{\rm SVEA}^{\max}\Phi_*)$。因此
\begin{align}
\frac{g(z',+\infty)}{g(z',-\infty)}
=\exp\Bigg\{
\ii\frac{q_{\rm e}}{\hbar\omega}
\operatorname{Im}\Bigg[
\ee^{\ii\kappa z'}
\int_{-\infty}^{+\infty}\dd z''\,
&E_{\omega,z}(z'')\ee^{-\ii\kappa z''}\\
&\times
 a_{\rm L}\!\left(\frac{z''-z'}{v_0}-\tau\right)
\Bigg]
+O(\epsilon_{\rm SVEA}^{\max}\Phi_*)
\Bigg\}.
\label{eq:pulse-phase-before-freeze}
\end{align}
因此在只使用 SVEA、尚未使用“穿越时间短”近似时，自然出现的是一个\emph{脉冲加权的轨迹 Fourier 投影}
\begin{equation}
\boxed{
\mathcal F_{\parallel}^{({\rm p})}(z',\tau)
\equiv
\int_{-\infty}^{+\infty}\dd z''\,
E_{\omega,z}(z'')\,
 a_{\rm L}\!\left(\frac{z''-z'}{v_0}-\tau\right)
\ee^{-\ii\kappa z''}.}
\label{eq:pulse-weighted-F}
\end{equation}
对应的慢变复耦合定义为
\begin{equation}
\boxed{
\beta_{\rm p}(z',\tau)
\equiv-\frac{q_{\rm e}}{2\hbar\omega}
\mathcal F_{\parallel}^{({\rm p})}(z',\tau).}
\label{eq:pulse-weighted-beta}
\end{equation}
于是式~\eqref{eq:pulse-phase-before-freeze} 可写成与连续波完全同构的形式
\begin{equation}
\boxed{
\frac{g(z',+\infty)}{g(z',-\infty)}
=
\exp\left[
-\beta_{\rm p}(z',\tau)\ee^{\ii\kappa z'}
+\beta_{\rm p}^*(z',\tau)\ee^{-\ii\kappa z'}
+O(\epsilon_{\rm SVEA}^{\max}\Phi_*)
\right].}
\label{eq:pulse-phase-beta-exact-svea}
\end{equation}
这条式子非常重要：若脉冲持续时间与电子穿越近场的时间可比，并不需要重新建立电子动力学，只需要保留式~\eqref{eq:pulse-weighted-F} 中空间场与时间包络的卷积；后面的“局域 $\beta$”只是它的进一步特殊极限。

现在才考察何时可以把光包络从近场积分中提出。设近场只在 $|z''|\lesssim L_{\rm nf}$ 内显著。对 Gaussian 包络，指数为
\begin{equation}
-\frac{(z''-z'-v_0\tau)^2}{4v_0^2\sigma_{\rm L}^2}.
\end{equation}
与 $z''=0$ 处相比，其变化量为
\begin{align}
 \left|
 \frac{(z''-z'-v_0\tau)^2-(-z'-v_0\tau)^2}
 {4v_0^2\sigma_{\rm L}^2}
 \right|
 &=\left|
 \frac{z''^2-2z''(z'+v_0\tau)}
 {4v_0^2\sigma_{\rm L}^2}
 \right|\ll1.
 \label{eq:pulse-freeze-exact}
\end{align}
对所有 $|z''|\le L_{\rm nf}$，一个充分条件是
\begin{equation}
 \boxed{
 \frac{L_{\rm nf}^2}{4v_0^2\sigma_{\rm L}^2}
 +\frac{L_{\rm nf}|z'+v_0\tau|}{2v_0^2\sigma_{\rm L}^2}
 \ll1.}
 \label{eq:pulse-freeze-condition}
\end{equation}
在电子与光脉冲主要重叠的区域 $|z'+v_0\tau|\lesssim v_0\sigma_{\rm L}$ 内，更直观的充分条件是
\begin{equation}
 \boxed{
 T_{\rm tr}\equiv\frac{L_{\rm nf}}{v_0}\ll\sigma_{\rm L}.}
 \label{eq:transit-condition}
\end{equation}
此时
\begin{equation}
 a_{\rm L}\!\left(\frac{z''-z'}{v_0}-\tau\right)
\simeq
 a_{\rm L}\!\left(-\frac{z'}{v_0}-\tau\right)
=
\exp\left[-\frac{(z'+v_0\tau)^2}{4v_0^2\sigma_{\rm L}^2}\right].
\label{eq:pulse-freeze}
\end{equation}
所以脉冲加权投影才进一步因式分解为
\begin{align}
\mathcal F_{\parallel}^{({\rm p})}(z',\tau)
&\simeq
\exp\left[-\frac{(z'+v_0\tau)^2}{4v_0^2\sigma_{\rm L}^2}\right]
\int_{-\infty}^{+\infty}
E_{\omega,z}(z'')\ee^{-\ii\kappa z''}\dd z''\\
&=
\exp\left[-\frac{(z'+v_0\tau)^2}{4v_0^2\sigma_{\rm L}^2}\right]
\Ftilde(\kappa),
\label{eq:pulse-F-factorization}
\end{align}
其中
\begin{equation}
 \kappa\equiv\frac{\omega}{v_0},
 \qquad
 \Ftilde(\kappa)
 \equiv\int_{-\infty}^{+\infty}
 E_{\omega,z}(z'')\ee^{-\ii\kappa z''}\,\dd z''.
 \label{eq:Fdef}
\end{equation}
因此
\begin{equation}
\beta_{\rm p}(z',\tau)
\simeq
\beta\,
\exp\left[-\frac{(z'+v_0\tau)^2}{4v_0^2\sigma_{\rm L}^2}\right],
\qquad
\beta=-\frac{q_{\rm e}}{2\hbar\omega}\Ftilde.
\label{eq:pulse-beta-factorization}
\end{equation}
代回式~\eqref{eq:pulse-phase-beta-exact-svea}，得到常用的局域相位形式
\begin{equation}
 \boxed{
 \frac{g(z',+\infty)}{g(z',-\infty)}
 =\exp\left\{
 \ii\frac{q_{\rm e}}{\hbar\omega}
 \operatorname{Im}\left[
 \ee^{\ii\kappa z'}\Ftilde(\kappa)
 \right]
 \exp\left[-\frac{(z'+v_0\tau)^2}{4v_0^2\sigma_{\rm L}^2}\right]
 \right\}.}
 \label{eq:phase-final}
\end{equation}
因此有限脉冲理论有清楚的两层结构：式~\eqref{eq:pulse-weighted-F} 是 SVEA 下仍保留穿越时间效应的一般结果；只有在 $T_{\rm tr}/\sigma_{\rm L}\ll1$ 后，才退化为“连续波 $\beta$ 乘局域光包络”的常用模型。


\section{从局域边带振幅到有限脉冲概率}

从本节开始额外采用式~\eqref{eq:transit-condition} 的短穿越条件，因此可以使用已经因式分解的局域耦合~\eqref{eq:pulse-beta-factorization}。若该条件不成立，应停留在脉冲加权耦合~\eqref{eq:pulse-weighted-beta}，而不能把连续波 $\beta$ 简单乘上到达时刻的光包络。有限脉冲概率仍不需要重新使用正、负能基底。记
\[
g_0(\zeta)\equiv g(\zeta,-\infty).
\]
耦合强度使用式~\eqref{eq:beta-standard} 定义的 $\beta$，并记 Bessel 自变量 $\eta=2|\beta|$：
\[
\eta=2|\beta|=\frac{e|\Ftilde|}{\hbar\omega}.
\]
对于高斯光脉冲，局域耦合参数为
\begin{equation}
\eta_{\rm loc}(\zeta,\tau)
=
\eta
\exp\left[
-\frac{(\zeta+v_0\tau)^2}
{4v_0^2\sigma_{\mathrm L}^2}
\right].
\label{eq:eta-local}
\end{equation}
电子电荷的负号只影响边带复振幅的相位约定；概率只依赖 $|\eta_{\rm loc}|$。

先说明有限脉冲概率为什么是“局域 Bessel 权重的积分”，而不是把这一形式当作经验平均。式~\eqref{eq:phase-final} 在固定 $\zeta=z'$ 处仍然只是一个正弦相位。把高斯实包络吸收到
\begin{equation}
\beta_{\rm loc}(\zeta,\tau)
\equiv
\beta\exp\!\left[-\frac{(\zeta+v_0\tau)^2}{4v_0^2\sigma_{\rm L}^2}\right],
\qquad
2|\beta_{\rm loc}|=\eta_{\rm loc},
\label{eq:beta-local-def}
\end{equation}
则正文已经固定的能量增益约定~\eqref{eq:sideband-amplitude-gain-convention} 可以逐点使用：
\begin{align}
g(\zeta,+\infty)
&=\sum_{n=-\infty}^{+\infty}
\ee^{+\ii n\kappa\zeta}\,g_n(\zeta),
\label{eq:pulse-sideband-decomposition}\\
g_n(\zeta)
&=g_0(\zeta)
J_n\!\left(2|\beta_{\rm loc}(\zeta,\tau)|\right)
\ee^{\ii n\arg[-\beta_{\rm loc}(\zeta,\tau)]}.
\label{eq:pulse-local-sideband-amplitude}
\end{align}
快速因子 $\ee^{+\ii n\kappa\zeta}$ 把第 $n$ 个分量的中心波数移到
$k_0+n\kappa$；$g_n(\zeta)$ 则是围绕该中心的慢包络。高斯光包络为实正数时，$\arg\beta_{\rm loc}=\arg\beta$，所以所有空间依赖都只进入 Bessel 自变量。
若相邻边带的中心波数间隔
\begin{equation}
\Delta k_{\rm sb}=\frac{\omega}{v_0}
\end{equation}
远大于单个边带自身的谱宽 $\Delta k_{\rm env}$，即
\begin{equation}
\boxed{\Delta k_{\rm env}\ll \omega/v_0,}
\label{eq:sideband-orthogonality-general}
\end{equation}
则不同 $n$ 的动量支撑近似互不重叠。第 $n$ 个谱峰的总概率可以单独积分。为避免把“中心波数平移”藏在记号里，先对第 $n$ 个完整分量
$h_n(\zeta)\equiv\ee^{\ii n\kappa\zeta}g_n(\zeta)$ 作 Fourier 变换。若
\begin{equation}
\widetilde g_n(k)=\frac{1}{\sqrt{2\pi}}
\int g_n(\zeta)\ee^{-\ii k\zeta}\dd\zeta,
\end{equation}
则逐项代入得
\begin{align}
\widetilde h_n(k)
&=\frac{1}{\sqrt{2\pi}}
\int g_n(\zeta)\ee^{-\ii(k-n\kappa)\zeta}\dd\zeta\\
&=\widetilde g_n(k-n\kappa).
\label{eq:finite-pulse-fourier-shift}
\end{align}
因此完整谱峰只是慢包络谱向 $k=n\kappa$ 的平移，积分范数不变。Parseval 定理给出
\begin{align}
P_n
&=\int_{-\infty}^{+\infty}|\widetilde g_n(k)|^2\dd k\\
&=\int_{-\infty}^{+\infty}|g_n(\zeta)|^2\dd\zeta\\
&=\int_{-\infty}^{+\infty}|g_0(\zeta)|^2
J_n^2\!\left(2|\beta_{\rm loc}(\zeta)|\right)\dd\zeta.
\label{eq:pulse-probability-parseval}
\end{align}
由于 $|\ee^{\ii n\arg(-\beta_{\rm loc})}|=1$，局域光学相位不会进入这个“整条边带积分概率”；但若探测器分辨边带内部的精细动量结构，或者不同边带互相重叠，就不能只保留式~\eqref{eq:pulse-probability-parseval}，而必须回到振幅级 Fourier 变换。

同一个公式也覆盖实验中的统计到达时间平均。由正文已经固定的关系 $t_c=-\zeta/v_0$，把电子穿过参考平面的到达时间记为
\begin{equation}
 t_{\rm e}\equiv t_c=-\frac{\zeta}{v_0}.
 \label{eq:arrival-time-zeta}
\end{equation}
若其归一化统计密度为

\begin{equation}
 \rho_{\rm e}(t_{\rm e})\ge0,
 \qquad
 \int_{-\infty}^{+\infty}\rho_{\rm e}(t_{\rm e})\dd t_{\rm e}=1,
 \label{eq:general-electron-time-density}
\end{equation}
并把局域复耦合写成
\begin{equation}
 \beta_{\rm loc}(t_{\rm e},\tau)
 =\beta_0\,a(t_{\rm e}-\tau),
 \label{eq:general-pulse-local-beta}
\end{equation}
其中 $a(t)$ 是归一到峰值 $|a|\le1$ 的光学慢包络。只要不同边带在动量空间可分辨，
\begin{equation}
 \boxed{
 P_n(\tau)=\int\dd t_{\rm e}\,
 \rho_{\rm e}(t_{\rm e})
 J_n^2\!\left(2|\beta_0a(t_{\rm e}-\tau)|\right).}
 \label{eq:general-pulse-convolution}
\end{equation}
所以有限脉冲问题本质上不是新的电子动力学，而是局域 Bessel 概率对电子到达时间分布的平均。

在弱耦合但允许观察任意固定阶 $|n|$ 时，Bessel 级数最低非零项为
\begin{equation}
 J_n(x)
 =\frac{1}{|n|!}\left(\frac{x}{2}\right)^{|n|}
 +O\!\left(x^{|n|+2}\right)
 \qquad (n\ge0\text{ 时无额外相位}),
 \label{eq:bessel-fixed-order-weak}
\end{equation}
故
\begin{equation}
 \boxed{
 P_n(\tau)
 \simeq
 \frac{|\beta_0|^{2|n|}}{(|n|!)^2}
 \int\dd t\,
 \rho_{\rm e}(t)
 |a(t-\tau)|^{2|n|}.}
 \label{eq:nth-order-cross-correlation}
\end{equation}
也就是说，第 $|n|$ 阶弱耦合延迟曲线是电子时间密度与“光强的 $|n|$ 次幂”的互相关。定义
\begin{equation}
 I_n(t)\equiv|a(t)|^{2|n|},
\end{equation}
则采用 Fourier 变换
\begin{equation}
 \widetilde f(\Omega)=\int\dd t\,f(t)\ee^{-\ii\Omega t}
\end{equation}
时，式~\eqref{eq:nth-order-cross-correlation} 的延迟 Fourier 谱可直接由 $u=t-\tau$ 求出：
\begin{align}
\widetilde P_n(\Omega)
&\propto\int\dd t\,\dd u\,
\rho_{\rm e}(t)I_n(u)
\ee^{-\ii\Omega(t-u)}\\
&=\widetilde\rho_{\rm e}(\Omega)
\widetilde I_n(-\Omega).
\end{align}
由于 $I_n(t)$ 为实函数，$\widetilde I_n(-\Omega)=\widetilde I_n^*(\Omega)$，因此
\begin{equation}
 \boxed{
 \widetilde P_n(\Omega)
 \propto
 \widetilde\rho_{\rm e}(\Omega)
 \widetilde I_n^{\,*}(\Omega).}
 \label{eq:delay-cross-correlation-fourier}
\end{equation}
因此实验上用高阶边带做时间门控，本质上是利用 $I_n(t)$ 随 $|n|$ 增大而变窄；Gaussian 情形中的 $1/\sqrt{|n|}$ 缩放只是这一一般互相关结构的特例。


\section{高斯脉冲：解析重叠、短/长脉冲与弱耦合极限}
若电子到达时间的标准差为 $\sigma_{\mathrm e}$，则纵向概率宽度为
\[
\sigma_z=v_0\sigma_{\mathrm e}.
\]
取归一化入射包络
\begin{equation}
g_0(\zeta)
=
\frac{1}{(2\pi\sigma_z^2)^{1/4}}
\exp\left(-\frac{\zeta^2}{4\sigma_z^2}\right),
\qquad
|g_0(\zeta)|^2
=
\frac{1}{\sqrt{2\pi}\sigma_z}
\exp\left(-\frac{\zeta^2}{2\sigma_z^2}\right).
\label{eq:gaussian-electron}
\end{equation}
当不同光子边带在动量空间不重叠，
\begin{equation}
\boxed{
\Delta k_{\rm env}\ll\frac{\omega}{v_0},}
\label{eq:special-nonoverlap}
\end{equation}
第 $n$ 个边带概率由 Parseval 定理写成
\begin{equation}
P_n(\tau)
=
\int_{-\infty}^{+\infty}
|g_0(\zeta)|^2
J_n^2\!\left[\eta_{\rm loc}(\zeta,\tau)\right]\dd\zeta.
\label{eq:Pn-general-gaussian}
\end{equation}
代入式 \eqref{eq:gaussian-electron} 与 \eqref{eq:eta-local}，
\begin{equation}
\boxed{
P_n(\tau)
=
\int_{-\infty}^{+\infty}
\frac{\dd\zeta}{\sqrt{2\pi}\,v_0\sigma_{\mathrm e}}
\exp\left[-\frac{\zeta^2}{2v_0^2\sigma_{\mathrm e}^2}\right]
J_n^2\!\left[
\eta
\exp\left(
-\frac{(\zeta+v_0\tau)^2}
{4v_0^2\sigma_{\mathrm L}^2}
\right)
\right].}
\label{eq:Pn-gaussian}
\end{equation}
这就是高斯电子脉冲与高斯光脉冲重叠时最直接的可观测边带概率。

上述高斯重叠式首先给出短电子脉冲极限。若电子包络在光脉冲变化尺度内足够窄，则可以把局域耦合在电子中心处展开。

若电子波包远短于光脉冲，
\begin{equation}
\boxed{\sigma_{\mathrm e}\ll\sigma_{\mathrm L},}
\label{eq:short-electron}
\end{equation}
则光学包络在电子概率分布支撑区内近似常数。更精确地，需要
\begin{equation}
\left|
\partial_\zeta\ln \eta_{\rm loc}(\zeta,\tau)
\right|_{\zeta\simeq0}
v_0\sigma_{\mathrm e}
=
\frac{|\tau|\sigma_{\mathrm e}}{2\sigma_{\mathrm L}^2}
\ll1,
\end{equation}
并且
\[
\frac{\sigma_{\mathrm e}^2}{\sigma_{\mathrm L}^2}\ll1.
\]
于是
\[
\eta_{\rm loc}(\zeta,\tau)\simeq \eta_{\rm loc}(0,\tau)
=
\eta\exp\left(-\frac{\tau^2}{4\sigma_{\mathrm L}^2}\right),
\]
归一化积分立即给出
\begin{equation}
\boxed{
P_n(\tau)
\simeq
J_n^2\!\left[
\eta
\exp\left(-\frac{\tau^2}{4\sigma_{\mathrm L}^2}\right)
\right].}
\label{eq:Pn-short-electron}
\end{equation}

同一公式的另一端是连续光或长光脉冲极限。此时光场包络在整个电子波包长度上几乎不变。

当
\begin{equation}
\boxed{
\sigma_{\mathrm L}\to\infty
\quad\text{或更一般地}\quad
\left|\partial_\zeta\ln\eta_{\rm loc}\right|L_{\mathrm e}\ll1,}
\label{eq:cw-condition}
\end{equation}
其中 $L_{\mathrm e}$ 是电子包络长度，局域耦合变为常数 $\eta_{\rm loc}\simeq\eta$。于是
\begin{equation}
\boxed{
P_n^{\rm CW}=J_n^2(\eta).}
\label{eq:Pn-cw}
\end{equation}
利用 Bessel 恒等式
\[
\sum_{n=-\infty}^{+\infty}J_n^2(\eta)=1
\]
可得
\[
\sum_nP_n^{\rm CW}=1.
\]
对于空间依赖耦合，只要边带仍满足式 \eqref{eq:special-nonoverlap}，同样有
\[
\sum_nP_n
=
\int |g_0(\zeta)|^2
\left[\sum_nJ_n^2(\eta_{\rm loc}(\zeta,\tau))\right]\dd\zeta
=1.
\]

最后，在上述任意脉宽关系之外还可以独立取弱耦合极限；它控制的是 Bessel 函数的幅度展开，而不是脉冲时宽。

若整个电子波包范围内
\begin{equation}
\boxed{|\eta_{\rm loc}(\zeta,\tau)|\ll1,}
\label{eq:weak-coupling-special}
\end{equation}
则
\begin{align}
J_0(\eta_{\rm loc})&=1-\frac{\eta_{\rm loc}^2}{4}+O(\eta_{\rm loc}^4),\\
J_{\pm1}(\eta_{\rm loc})&=\pm\frac{\eta_{\rm loc}}{2}+O(\eta_{\rm loc}^3),\\
J_{|n|\ge2}(\eta_{\rm loc})&=O(\eta_{\rm loc}^{|n|}).
\end{align}
于是
\begin{align}
P_{\pm1}(\tau)
&\simeq
\frac14
\int |g_0(\zeta)|^2\eta_{\rm loc}^2(\zeta,\tau)\dd\zeta,\\
P_0(\tau)
&\simeq
1-\frac12
\int |g_0(\zeta)|^2\eta_{\rm loc}^2(\zeta,\tau)\dd\zeta.
\label{eq:weak-prob-general}
\end{align}
对式 \eqref{eq:gaussian-electron}，令
\[
s=\frac{\zeta}{v_0},
\qquad
\dd\zeta=v_0\dd s,
\]
则
\begin{align}
\int |g_0|^2\eta_{\rm loc}^2\,\dd\zeta
&=
\eta^2
\int_{-\infty}^{+\infty}
\frac{\dd s}{\sqrt{2\pi}\sigma_{\mathrm e}}
\exp\left[
-\frac{s^2}{2\sigma_{\mathrm e}^2}
-\frac{(s+\tau)^2}{2\sigma_{\mathrm L}^2}
\right].
\end{align}
把指数配方。令
\[
A=\frac{1}{\sigma_{\mathrm e}^2}
+\frac{1}{\sigma_{\mathrm L}^2},
\qquad
B=\frac{\tau}{\sigma_{\mathrm L}^2},
\]
则
\begin{align}
\frac{s^2}{\sigma_{\mathrm e}^2}
+\frac{(s+\tau)^2}{\sigma_{\mathrm L}^2}
&=
As^2+2Bs+\frac{\tau^2}{\sigma_{\mathrm L}^2}\\
&=
A\left(s+\frac{B}{A}\right)^2
+\frac{\tau^2}{\sigma_{\mathrm e}^2+\sigma_{\mathrm L}^2}.
\end{align}
因此
\begin{align}
\int |g_0|^2\eta_{\rm loc}^2\,\dd\zeta
&=
\eta^2
\frac{1}{\sqrt{2\pi}\sigma_{\mathrm e}}
\exp\left[
-\frac{\tau^2}{2(\sigma_{\mathrm e}^2+\sigma_{\mathrm L}^2)}
\right]
\int_{-\infty}^{+\infty}
\exp\left[
-\frac{A}{2}\left(s+\frac{B}{A}\right)^2
\right]\dd s\\
&=
\eta^2
\frac{\sigma_{\mathrm L}}
{\sqrt{\sigma_{\mathrm e}^2+\sigma_{\mathrm L}^2}}
\exp\left[
-\frac{\tau^2}
{2(\sigma_{\mathrm e}^2+\sigma_{\mathrm L}^2)}
\right].
\label{eq:gaussian-overlap}
\end{align}
于是弱耦合下
\begin{equation}
\boxed{
P_{\pm1}(\tau)
\simeq
\frac{\eta^2}{4}
\frac{\sigma_{\mathrm L}}
{\sqrt{\sigma_{\mathrm e}^2+\sigma_{\mathrm L}^2}}
\exp\left[
-\frac{\tau^2}
{2(\sigma_{\mathrm e}^2+\sigma_{\mathrm L}^2)}
\right],}
\label{eq:P1-gaussian-weak}
\end{equation}
并且
\begin{equation}
\boxed{
P_0(\tau)
\simeq
1-
\frac{\eta^2}{2}
\frac{\sigma_{\mathrm L}}
{\sqrt{\sigma_{\mathrm e}^2+\sigma_{\mathrm L}^2}}
\exp\left[
-\frac{\tau^2}
{2(\sigma_{\mathrm e}^2+\sigma_{\mathrm L}^2)}
\right].}
\label{eq:P0-gaussian-weak}
\end{equation}
因此弱耦合下的 PINEM 延迟曲线就是电子与光脉冲强度包络的高斯互相关，其方差为
\[
\sigma_{\rm cross}^2
=\sigma_{\mathrm e}^2+\sigma_{\mathrm L}^2.
\]


\section{把局域耦合重新组装成统一有限脉冲边带}

将前面的局域耦合、有限脉冲和边带展开合并，可把这一近似层级的结果写成统一形式。
在一阶、准单色、弱反冲和边带不重叠条件同时成立时，PINEM 的核心可压缩为
\begin{equation}
\boxed{
g_{\rm out}(\zeta)
=
g_0(\zeta)
\exp\left\{
\ii\frac{q_{\rm e}}{\hbar\omega}
\operatorname{Im}\left[
\ee^{\ii(\omega/v_0)\zeta}
\Ftilde\!\left(\frac{\omega}{v_0}\right)
\right]
G(\zeta,\tau)
\right\},}
\label{eq:unified-master-phase}
\end{equation}
其中
\[
G(\zeta,\tau)
=
\exp\left[
-\frac{(\zeta+v_0\tau)^2}
{4v_0^2\sigma_{\mathrm L}^2}
\right].
\]
Jacobi--Anger 展开给出
\begin{equation}
\boxed{
g_{\rm out}(\zeta)
=
\sum_{n=-\infty}^{+\infty}
g_n(\zeta)
\ee^{\ii n(\omega/v_0)\zeta},}
\end{equation}
并且
\begin{equation}
\boxed{
g_n(\zeta)
=
g_0(\zeta)
\left(-\frac{\Ftilde}{|\Ftilde|}\right)^n
J_n\left[
-\frac{q_{\rm e}|\Ftilde|}{\hbar\omega}
G(\zeta,\tau)
\right].}
\label{eq:unified-gn}
\end{equation}
式~\eqref{eq:unified-master-phase} 已经给出有限脉冲下的完整标量边带包络。若还需要恢复每个通道的 Dirac 自旋量及负能小分量修正，可在附录~\ref{app:dirac-two-band-check} 中从同一正能包络逐项重构；这些修正不参与本节的边带概率与时间平均。
